\documentclass[twoside,11pt,openright]{report}

\usepackage[latin1]{inputenc}
\usepackage[american]{babel}
\usepackage{a4}
\usepackage{latexsym}
\usepackage{amssymb}
\usepackage[fleqn]{amsmath}
\usepackage{epsfig}
\usepackage[T1]{fontenc}
\usepackage{mathptmx}
\usepackage{color}
\usepackage{epstopdf}
\usepackage{microtype}
\usepackage{hyperref}
\usepackage[useregional]{datetime2}
\DTMlangsetup[en-US]{showdayofmonth=false}
\usepackage{lipsum}

\usepackage{amsthm}
\theoremstyle{definition}
\newtheorem{definition}{Definition}[section]
\newcommand{\definitionautorefname}{Definition}


\theoremstyle{plain}
\newtheorem{lemma}{Lemma}[section]

\renewcommand*\sfdefault{lmss}
\renewcommand*\ttdefault{txtt}

\newcommand{\todo}[1]{{\color[rgb]{.5,0,0}\textbf{$\blacktriangleright$#1$\blacktriangleleft$}}}

\begin{document}

%%%%%%%%%%%%%%%%%%%%%%%%%%%%%%%%%%%%%%%%%%%%%%%%%%%%%%%%%%%%%%%%%%%%%%%

\pagestyle{empty}
\pagenumbering{roman}
\vspace*{\fill}\noindent{\rule{\linewidth}{1mm}\\[4ex]
{\Huge\sf SD-in-the-Head rust implementation and optimization}\\[2ex]
{\huge\sf Hugh Benjamin Zachariae, 201508592 \\ Magnus Jensen,
201708626}\\[2ex]
\noindent\rule{\linewidth}{1mm}\\[4ex]
\noindent{\Large\sf Master's Thesis, Computer Science\\[1ex]
  \today \\[1ex] Advisor: Diego F. Aranha\\[15ex]}\\[\fill]}
\epsfig{file=logo.eps}\clearpage

%%%%%%%%%%%%%%%%%%%%%%%%%%%%%%%%%%%%%%%%%%%%%%%%%%%%%%%%%%%%%%%%%%%%%%%

\pagestyle{plain}
\chapter*{Abstract}
\addcontentsline{toc}{chapter}{Abstract}

\todo{in English\dots}

\chapter*{Resum\'e}
\addcontentsline{toc}{chapter}{Resum\'e}

\todo{in Danish\dots}

\chapter*{Acknowledgments}
\addcontentsline{toc}{chapter}{Acknowledgments}

\todo{\dots}

\vspace{2ex}
\begin{flushright}
  \emph{Magnus Jensen,}\\
  \emph{Aarhus, \today.}
\end{flushright}

\tableofcontents
\cleardoublepage
\pagenumbering{arabic}
\setcounter{secnumdepth}{2}

%%%%%%%%%%%%%%%%%%%%%%%%%%%%%%%%%%%%%%%%%%%%%%%%%%%%%%%%%%%%%%%%%%%%%%%

\chapter{Introduction}
\label{ch:intro}

few pages. INtroduce what we have done and how the paper is structured

\todo{motivate and explain the problem to be addressed}

\todo{example of a citation: \cite{DBLP:conf/sas/ChristensenMS03}}
\todo{get your bibtex entries from \url{https://dblp.org/}}


\section{Quantum Security}
\label{sec:nist}

In 1994, the discovery of Shor's algorithm \cite{shor1997}, provided a polynomial time algorithm for factoring integers and computing discrete logarithms using quantum computing. Back then, quantum computers where theories on paper. However, two decades later, significant developments in quantum technology, has brought us dangerously close to the point where quantum computers can solve foundational problems, which would break the security of modern standards of public-key cryptosystems. For example the Diffie-Hellman key exchange standard X25519, used in TLSv1.3, is not post-quantum secure.

In 2016, the National Institute of Standards and Technology (NIST), known for standardizing cryptographic primitives like AES and SHA, initiated a call for quantum resistant public-key cryptosystems \cite{nistcall}. This let to the standardization of three signature schemes, namely Dilithium~\cite{ducas2018crystals}, Falcon~\cite{fouque2018falcon} and SPHINCS+~\cite{bernstein2019sphincs+}. All three of the schemes are based on the hardness of structured lattice problems. This has led to the most recent call in 2022 for additional quantum resistant signature schemes based on different assumptions to diversify the standards available for the future.


\subsection{NIST security evaluation}

With the rapid development in quantum computing, there remains

AES-128, AES-192, AES-256

%%%%%%%%%%%%%%%%%%%%%%%%%%%%%%%%%%%%%%%%%%%%%%%%%%%%%%%%%%%%%%%%%%%%%%%

\chapter{Description of the algorithm}
\label{ch:desc}

Abstract level description of the algorithm own words
MPCitH
Syndrome detection problem
Fiat-Shamir heuristic




\section{Syndrome Decoding Problem}
\label{sec:syndrome}

Describe the the syndrome decoding problem introducing why we need ZKAoK, LSSS, Beaver triples and MPC construction from \cite{baum2020concretely} section 3.3.

\begin{definition}
  \label{def:syndrome}
  Given $H \in \mathbb{F}^{(m-k)\times m}_q$ and $y \in \mathbb{F}^{m-k}_q$. The problem is to find $x \in \mathbb{F}^m_q$ s.t. wt$(x) = w$ and $Hx = y$.
\end{definition}


\section{Galois Finite Field}
\label{sec:gf256}

A group for bytes

\todo{Describe the arithmetic operations in GF256 and the extension field}

\subsection{Extension field}
$F_q^\eta$

\subsubsection{Polynomial evaluation}
For polynomial evalution in the multiparty computation seen in \autoref{sec:mpc}, we need to evaluate a polynomial $P(x)$ at a point $r \in F_q^\eta$. This is done by evaluating the polynomial at each coefficient and summing the results. The polynomial $P(x)$ is defined as
\begin{align}
  \textstyle\bigcup_{|Q|}(F_q)^{|Q|} \times F_q^\eta & \rightarrow F_q^\eta                 \nonumber                                  \\
  Q(r)                                               & = \textstyle\sum_{i=0}^{|Q|} Q_i \cdot r^{i-1} &  & Q_i \in F_q, r \in F_q^\eta
  \label{eq:mpcpoly}
\end{align}

\section{Linear Secret Sharing Scheme (LSSS)}
\label{sec:sss}

\textit{Secret sharing schemes} (SSS) are a type of cryptographic protocol that allows for the distribution of a secret amongst a group of participants. The secret can only be reconstructed when a sufficient number of shares are combined together. The threshold variant of the SDitH protocol relies on a low-threshold linear secret sharing scheme (\textit{LSSS}). Threshold secret sharing schemes allow for the reconstruction of a secret from a subset of shares of length $\ell$, where $\ell$ is the threshold. The threshold allows for the SDitH protocol to be more communication efficient, as the amount of shares needed to reconstruct the secret is low.

\begin{definition}
  \label{def:sss}
  \textit{$S$ is a $(\ell,n)$ threshold SSS if it satisfies the following properties}:

  \begin{itemize}
    \item \textbf{Share generation}: Given a secret $s$, the scheme generates $n$ shares $\texttt{share}(s) = [[s]] = [[s_1, s_2, \dots, s_n]]$.
    \item \textbf{Reconstruction}: Given a subset of shares $[[s']]$ of size $\ell$, the scheme can reconstruct the secret $s = \texttt{open}([[s']])$.
  \end{itemize}

\end{definition}

\textit{Linear secret charing schemes}, or $(+,+)$-homomorphic schemes, refers to secret sharing schemes that are linearly homomorphic over some field $\mathbb{F}$ (say the galois field \texttt{GF256} described in \autoref{sec:gf256}). This means that given shares $[[a]]$ and $[[b]]$ we have that
\begin{definition}
  A $(\ell,n)$ threshold SSS is $(+,+)$-homomorphic if for any two secrets $s_1$ and $s_2$ and their shares $[[s_1]]$ and $[[s_2]]$, the sum of the shares $[[s_1]] + [[s_2]] = [[s_{11} + s_{21}, s_{12} + s_{22}, \dots, s_{1n} + s_{2n}]]$ is equal to the share of the sum of the secrets $[[s_1 + s_2]]$ for the same subset of shares.
\end{definition}
\subsection{Shamir's Secret Sharing}
\label{sec:shamir}

Shamir's Secret Sharing scheme \cite{shamir1979share}. Shamir's Secret Sharing is a method for distributing a secret amongst a group of participants, each of which is given a share of the secret. The secret can only be reconstructed when a sufficient number of shares are combined together.

For a secret $s$ and a given security threshold $\ell$, the share $[[s]]$ is generated by sampling a random polynomial of degree $t-1$ with $s$ as the free coefficient. Each participant is given a share of the polynomial evaluated at a point. The secret can be reconstructed by interpolating the polynomial from a sufficient number of shares.

The $(t,n)$ Shamir's polynomial based secret sharing scheme is $(+,+)$-homomorphic in which the addition of two polynomials secrets equals the Lagrange's interpolation of the sum-of-shares for the same subset of shares.

\todo{Write about creating shares, addition}

\subsection{Multiplication and Beaver Triples}
\label{sec:beaver}

\section{Zero-Knowledge Proofs}
\label{sec:zkp}

Zero-Knowledge Arguments of Knowledge (ZKAoK) allow a prover to convince a verifier that they knows a certain witness for a statement without revealing any additional information about the witness. \todo{What are they used for..}

ZKAoK protocols have been introduced in multiple variants that each solve specific relational problems with varying levels of communication complexity and prover running times. While many of these perform well for these specific problems, there has long been a need for a well performing generic ZKAoK protocol.

\subsection{MPC-in-the-Head}
\label{sec:mpcith}

Secure Multiparty Computation or MPC describes a cryptographic protocol that allows multiple parties to run a circuit $C$, while no information on inputs is leaked outside of the intended output of $C$. These algorithms can be augmented to ensure correctness, even in the presence of malicious parties. \todo{skriv mere her. Check threshold paper måske?}

The MPC-in-the-Head (MPCitH) framework, introduced by~\cite{ishai2007zero}, uses MPC protocols to construct generic ZKAoK protocols. First, the statement is turned into a circuit $C$. On input $x$, $C$ outputs $y$ iff. $x = w$ for some witness $w$ which is a correct witness for the statement to be proven. The prover simulates the MPC protocol of all the parties \textit{in the head}. The parties each evalutate $C$ on a secret sharing of $w$

The framework benefits from previous efficiency improvements in MPC protocols to improve the efficiency of ZKAoK protocols. Particularly, its possible to achieve \textit{constant-rate} zero knowledge. For an arbitrary circuit $C$ of size $s$ and a bounded fan-in, it is possible to construct a ZKAoK protocol with communication complexity $O(s) + \text{poly}(k)$, for a security parameter $k$. This is an improvement over previous ZKAoK protocols, which had communication complexity $O(ks)$. \todo{write about this in the previous section}.

\subsection{MPC preprocessing}

We first generalize the idea of [KKW18] to work over arithmetic circuits using a variant of the SPDZ MPC protocol [DPSZ12,LN17] and provide a formal proof of security to their "cut-and-choose" preprocessing heuristic. Then, we present a new construction where we replace the "cut-and-choose" mechanism with a "sacrificing"-based approach. While both approaches have similar cost per MPC instance, our "sacrificing"-based approach yields a smaller cheating probability, which means that the number of MPC instances simulated in the proof can be significantly smaller, thus reducing the overall communication footprint. Our scheme is highly flexible in its choice of parameters. In particular, by changing the number of parties in the underlying MPC protocol, one can alternate between achieving low communication and low running time. Our construction only requires efficient standard symmetric primitives, and thus is plausibly post-quantum secure even in the non-interactive case [DFMS19]. The two constructions can be found in \cite{baum2020concretely} \todo{Copied from citation. rewrite this}

\subsection{Verification of a multiplication triple using another}
\label{sec:verify}
Given a random triple $([[a]], [[b]], [[c]])$, it is possible to verify the correctness of a triple $([[x]], [[y]], [[z]])$, i.e., that $z = x \cdot y$, without revealing any information on either of the triples, in the following way \cite{baum2020concretely}.

\begin{enumerate}
  \item The parties generate a random $\epsilon \in \mathbb{F}$.
  \item The parties locally set $[[\alpha]] = \epsilon[[x]] + [[a]], [[\beta]] = [[y]] + [[b]]$.
  \item The parties run \texttt{open}$([[\alpha]])$ and \texttt{open}$([[\beta]])$ to obtain $\alpha$ and $\beta$.
  \item The parties locally set $[[v]] = \epsilon[[z]] - [[c]] + \alpha  \cdot [[b]] + \beta  \cdot [[a]] - \alpha  \cdot \beta$.
  \item The parties run \texttt{open}$([[v]])$ to obtain $v$ and accept iff $v = 0$.

\end{enumerate}
Observe that if both triples are correct multiplication triples (i.e., $z = xy$ and $c = ab$) then the parties will always accept since
\begin{align}
  v & = \epsilon \cdot z - c + \alpha \cdot b + \beta \cdot a - \alpha \cdot \beta                                            \\
    & = \epsilon \cdot xy - ab + (\epsilon \cdot x + a)b + (y + b)a - (\epsilon \cdot x + a)(y + b)                           \\
    & = \epsilon \cdot xy - ab + \epsilon \cdot xb + ab + ya + ba - \epsilon \cdot xy - \epsilon \cdot xb - ay - ab           \\
    & = (\epsilon \cdot xy - \epsilon \cdot xy) + (ab - ab) + (\epsilon \cdot xb - \epsilon \cdot xb) + (ya - ay) + (ba - ab) \\
    & = 0
\end{align}

\subsection{Error probability}

\begin{lemma}
  If $([[a]], [[b]], [[c]])$ or $([[x]], [[y]], [[z]])$ is an incorrect multiplication triple then the parties output \texttt{accept} in the sub-protocol above with probability $\frac{1}{|\mathbb{F}|}$.
\end{lemma}

\todo{Proof of lemma and solution}

\section{Fiat-Shamir Heuristic}
\label{sec:fiatshamir}

\todo{Turning the ZKAoK into a signature scheme}

%%%%%%%%%%%%%%%%%%%%%%%%%%%%%%%%%%%%%%%%%%%%%%%%%%%%%%%%%%%%%%%%%%%%%%%

\chapter{Specification}
\label{ch:spec}
more detailed description of the algorithm
e.g. how we sampled I[e] witness challenge
table of spec params (with our code naming and different categories)

\section{Merkle Tree Commitment}

\section{Hashing and XOF}

\todo{Write about Shake, Keccak} and how we initiate the hash function and XOF

\section{MPC computation}

The computation is based on HVZKAoK Protocol using imperfect preprocessing and sacrificing. Section 3.3 \cite{baum2020concretely}.

A toy example of the computation protocol computations, can be seen in~\autoref{fig:mpc}. We have the two computation methods. Here for 1 split and 1 evaluation point. This means that we have only value for each challenge and beaver triple. Note that any arithmetic is run in GF256, so addition and subtraction are both \texttt{XOR} and multiplication is modulus $x^8 + x^4 + x^3 + x + 1$. Furthermore, for negation we have that $-a = a$.

\begin{figure}
  \makebox[\textwidth]{
    \fbox{\begin{minipage}[t]{.45\textwidth}
        \noindent \texttt{PartyComputation}
        \begin{flalign*}
           & \textit{Input: }                                                               \\
           & (s_a, Q', P, a, b, c), (\overline{\alpha}, \overline{\beta}), (H', y)          \\
           & (\epsilon, r), \texttt{with\_offset}                                           \\
           & \textit{Output: }                                                              \\
           & (\alpha, \beta, v)                                                             \\\\
           & Q = Q'_1\text{ if \texttt{with\_offset} else }Q'_0                             \\
           & S = (s_a | y + H's_a) \text{ if \texttt{with\_offset} else } (s_a | H's_a)     \\
           & v = -c                                                                         \\
           & \alpha = \epsilon \cdot Q(r) + a                                               \\
           & \beta = S(r) + b                                                               \\
           & v \mathrel{{+}{=}} \epsilon \cdot F(r) \cdot P(r)                              \\
           & v \mathrel{{+}{=}} \overline{\alpha} \cdot b + \overline{\beta} \cdot a        \\
           & v \mathrel{{+}{=}} - \alpha \cdot \beta \text{ \ \ \ if \texttt{with\_offset}}
        \end{flalign*}
      \end{minipage}}
    \hfill
    \noindent
    \fbox{\begin{minipage}[t]{.45\textwidth}
        \noindent \texttt{InverseComputation}
        \begin{flalign*}
           & \textit{Input: }                                                                 \\
           & (s_a, Q', P), (\alpha, \beta, v), (\overline{\alpha}, \overline{\beta}), (H', y) \\
           & (\epsilon, r), \texttt{with\_offset}                                             \\
           & \textit{Output: }                                                                \\
           & (a, b, c)                                                                        \\\\
           & Q = Q_1\text{ if \texttt{with\_offset} else }Q_0                                 \\
           & S = (s_a | y + H's_a) \text{ if \texttt{with\_offset} else } (s_a | H's_a)       \\
           & c = -v                                                                           \\
           & a = \alpha - \epsilon \cdot Q(r)                                                 \\
           & b = \beta - S(r)                                                                 \\
           & c \mathrel{{+}{=}} \epsilon \cdot F(r) \cdot P(r)                                \\
           & c \mathrel{{+}{=}} \overline{\alpha} \cdot b + \overline{\beta} \cdot a          \\
           & c \mathrel{{+}{=}} - \alpha \cdot \beta \text{ \ \ \ if \texttt{with\_offset}}
        \end{flalign*}
      \end{minipage}}}
  \caption{Simplified version of the MPC party computation and inverse computation. $Q_0$ means that $Q$ is completed with a $0$ for leading coefficient. Furthermore, $F$ is precomputed. Note that all arithmetic is done in $\mathbb{F}_q = GF256$. All elements are in $\mathbb{F}_q^\eta$ except for the coefficients of $Q$, $S$ and $P$ which are in $\mathbb{F}_q$.}
  \label{fig:mpc}
\end{figure}
\bigskip

Note that the

If we first instantiate an input $i$ and one random input $i^*$ (like the \texttt{input\_coef}).Then the input share is generated by adding the two. Similar, but simpler, to the input share generation of Algorithm 12, line 13 of the specification.
\begin{align*}
  i             & = (s_a, Q, P, a,b,c)                                             \\
  i^*           & = ({s_a}^*, Q^*, P^*, a^*,b^*,c^*)                               \\
  [i] = i + i^* & = ({s_a} + {s_a}^*, Q + Q^*, P + P^*, a + a^*, b + b^*, c + c^*) \\
                & = ([s_a], [Q], [P], [a], [b], [c])                               \\
  \texttt{chal} & = (\epsilon, r)                                                  \\
  \texttt{pk}   & = (H', y)
\end{align*}
We also compute the plain broadcast share of the input as per Algorithm 12, line 18. Note that $\overline{v}$ is computed to zero and therefore removed from the computation in the implementation.
\begin{align*}
  (\overline{\alpha},\ \overline{\beta})         & =
  \texttt{PartyComputation}(i,\ (\overline{\alpha}, \overline{\beta}),\ \texttt{chal},\ \texttt{pk},\ \texttt{true})                                                                                   \\\\
  \overline{\alpha}                         =    & \epsilon \cdot Q_1(r) + a                                                                                                                           \\
  \overline{\beta}                          =    & S_y(r)  + b                                                                                                                                         \\
  \overline{v}                            =      & -c + \epsilon \cdot F(r) \cdot P(r) + \overline{\alpha} \cdot b + \overline{\beta} \cdot a  - \overline{\alpha} \cdot \overline{\beta}              \\
  \overline{v}                            =      & -c + \epsilon \cdot F(r) \cdot P(r) + (\epsilon \cdot Q_1(r) + a) \cdot b + (S_y(r)  + b) \cdot a  - (\epsilon \cdot Q_1(r) + a) \cdot (S_y(r) + b) \\
  \overline{v}                            =      & -c + \epsilon \cdot F(r) \cdot P(r)                                                                                                                 \\
                                                 & + \epsilon \cdot Q_1(r) \cdot b + c + S_y(r) \cdot a + c                                                                                            \\
                                                 & - \epsilon \cdot Q_1(r) \cdot S_y(r) - \epsilon \cdot Q_1(r) \cdot b - a \cdot S_y(r) - c                                                           \\
  \overline{v}                                 = & \ 0
\end{align*}
We then compute a broadcast share from the randomness and the broadcast, as per Algorithm 12, line 21.
\begin{align*}
  (\alpha^*, \beta^*, v^*) & = \texttt{PartyComputation}(i^*, (\overline{\alpha},
  \overline{\beta}), \texttt{chal}, \texttt{pk}, \texttt{false})                  \\\\
  \alpha^*                 & = \epsilon \cdot {Q^*}_0(r) + a^*                    \\
  \beta^*                  & =  {S^*}_0(r) + b^*
\end{align*}
This broadcast share is sent to the verifier along with the truncated input share (removing the beaver triples). The verifier then needs to recompute the input share beaver triples using the \texttt{InverseComputation} function. First we add the input share to the broadcast share as per Algorithm 13, line 8.
\begin{align*}
  (\alpha', \beta', v') & = (\alpha^*, \beta^*, v^*) + (\overline{\alpha},\ \overline{\beta}, 0)
  = (\alpha^* + \overline{\alpha}, \beta^* + \overline{\beta}, v^* + 0)                          \\\\
  \alpha'               & = \epsilon \cdot {Q^*}_0(r) + a^* + \overline{\alpha}                  \\
                        & = \epsilon \cdot {Q^*}_0(r) + a^* + \epsilon \cdot Q_1(r) + a          \\
  \beta'                & = {S^*}_0(r) + b^* + \overline{\beta}                                  \\
                        & = {S^*}_0(r) + b^* + S_y(r) + b                                        \\
\end{align*}
Next, the verifier computes the inverse of the broadcast share to recompute $([a], [b], [c])$ using the \texttt{InverseComputation} function. This is done as per Algorithm 13, line 10.
\begin{align*}
  (a', b', c') & = \texttt{InverseComputation}([i], (\alpha', \beta', v'),
  (\overline{\alpha}, \overline{\beta}), \texttt{chal}, \texttt{pk}, \texttt{true})                                                  \\\\
  a'           & = \alpha' - \epsilon \cdot {[Q]}_1(r)                                                                               \\
               & = \epsilon \cdot {Q^*}_0(r) + a^* + \epsilon \cdot Q_1(r) + a - \epsilon \cdot {[Q]}_1(r)                           \\
               & = \epsilon \cdot {Q^*}_0(r) - \epsilon \cdot {[Q]}_1(r) + \epsilon \cdot Q_1(r) + [a]                               \\
               & = \epsilon \cdot ({Q^*}_0(r) - {[Q]}_1(r) +  Q_1(r)) + [a]                                                          \\
               & = \epsilon \cdot (Q^*_0(r) +  Q^*_0(r)) + [a]                                             &  & \autoref{eq:mpcpoly} \\
               & = [a]                                                                                     &  & \autoref{eq:mpcpoly} \\
  b'           & = \beta' - {[S]}_y(r)                                                                                               \\
               & = {S^*}_0(r) + b^* + S_y(r) + b - {[S]}_y(r)                                                                        \\
               & = {S^*}_0(r) - {[S]}_y(r) + S_y(r)  + [b]                                                                           \\
               & = {S^*}_0(r) - S^*_0(r)  + [b]                                                            &  & \autoref{eq:mpcpoly} \\
               & = [b]                                                                                     &  & \autoref{eq:mpcpoly} \\
\end{align*}
\todo{Want to explain the above in a more detailed manner? Specifically ${Q^*}_0(r) - {[Q]}_1(r) = Q_1(r)$}

\section{Security}

EUF-CMA: unforgability against chosen message attacks

Random $pk$, and ask signature oracle to sign messages $(m_1, ..., m_r)$, then forge $m, \sigma$ s.t. $\sigma$ is a valid signature for $m$ for $pk$

\subsection{Assumptions}

Syndrome Decoding instances cannot be solved in complexity lower than $2^\kappa$ corresponding to the complexity of breaking AES by exhaustive search. For this $\kappa$ is defined to be 128 for category 1 (AES-128), 192 for category 3 (AES-192) and 256 for category 5 (AES-256).


XOF is secure with 128-bit, 192-bit, 256-bit security levels.

Hash functions behave as random oracles. Security holds in Random Oracle Model (ROM) and Quantum Random Oracle Model (QROM).

\subsection{Security of the Syndrome Detection Problem}
\label{sec:sdsec}

Recall the definition of the syndrome detection problem from \autoref{def:syndrome}. A brute force attack of guessing $x$ would require finding a single correct solution in $\binom{m}{w} q^w$. This would be infeasible for the parameters.

Best known algorithms: Generalized Birthday Algorithms (GBA) and Information Set Decoding (ISD)

\subsection{Security parameters}
\label{sec:secparam}

\subsubsection{SD Parameters}
\label{sec:sdparam}

The parameters for the SD problem are chosen according to the estimated cost for the best known algorithms for solving such problems. GBA and ISD described in \autoref{sec:sdsec}.

\subsubsection{MPC Parameters}

\subsubsection{Cryptographic primitives}


%%%%%%%%%%%%%%%%%%%%%%%%%%%%%%%%%%%%%%%%%%%%%%%%%%%%%%%%%%%%%%%%%%%%%%%

\chapter{Implementation}
\label{ch:impl}

Tooling and language feaures (rust, criterion)
code sections
code re-usability with traits for categories.

const generics vs inline mutability (benchmarking?, nightly?)
traits for categories

%%%%%%%%%%%%%%%%%%%%%%%%%%%%%%%%%%%%%%%%%%%%%%%%%%%%%%%%%%%%%%%%%%%%%%%

\chapter{Benchmarks}
\label{ch:bench}
diaries of benchmarks.
discussion of results

%%%%%%%%%%%%%%%%%%%%%%%%%%%%%%%%%%%%%%%%%%%%%%%%%%%%%%%%%%%%%%%%%%%%%%%

\chapter{Conclusion}
\label{ch:conclusion}

wrap up and pose future work
what should people continue with
point to round 2 NIST
work in context of timeline

\todo{conclude on the problem statement from the introduction}

%%%%%%%%%%%%%%%%%%%%%%%%%%%%%%%%%%%%%%%%%%%%%%%%%%%%%%%%%%%%%%%%%%%%%%%

\cleardoublepage
\addcontentsline{toc}{chapter}{Bibliography}
\bibliographystyle{plain}
\bibliography{refs}

%%%%%%%%%%%%%%%%%%%%%%%%%%%%%%%%%%%%%%%%%%%%%%%%%%%%%%%%%%%%%%%%%%%%%%%

\cleardoublepage
\appendix
\chapter{The Technical Details}

\todo{\dots}

\end{document}
